% !TEX root = ../main.tex
\chapter{서론 및 연구 문제}
\label{ch:introduction}

탈중앙화 금융(DeFi)은 전통적 금융 중개자 없이 대출, 차입, 레버리지 거래를 가능하게 함으로써 공개 블록체인에서 자본이 배분되는 방식을 변화시켰다 \cend{53}. 그러나 대부분의 신용형 상호작용은 여전히 높은 담보를 요구하는데, 이는 프로토콜이 오프체인 신원 체계나 전통적 신용 이력에 의존할 수 없기 때문이다 \cend{48}. 가명 환경에서 운영상의 과제는 자본을 안전하게 배분하는 것, 즉 레버리지 노출을 신뢰성 있게 관리하는 지갑과 반복적 손실, 청산, 적대적 행동을 보이는 지갑을 구분하는 것이다.

본 학위논문은 그 설정 안의 구체적 연구 문제를 다룬다. ERC-20 \texttt{approve}/\texttt{permit} allowance로부터 구축한 온체인 평판 점수---\textbf{EndorseRank}로 구현---가 계산 비용을 줄이고, 보증(endorsement) 의미를 더 충실히 반영하며, \textbf{Arbitrum One}에서 전송 기반 Adaptive Weighted PageRank(AWP)와 비교할 때 구별되면서도 신뢰할 만한 정합 프로파일을 산출할 수 있는가이다. 주요 실증 비교는 동일 지갑 표본에서 \textbf{EndorseRank 대 AWP}이며, Do et al.\cend{16}의 순위 상관 틀 아래 다섯 가지 대리 지표 계열(GMX V2 무기한 결과 포함)로 검증한다.

본 연구에서 \textbf{온체인 평판 점수} 계산이란 오프체인 신원 없이 관측 가능한 온체인 행동과 관계로부터 정량적 지갑 점수를 도출하는 것을 말한다 \cend{16}. EndorseRank는 allowance 기반 보증 간선으로부터 그러한 점수를 산출한다. 본 장의 나머지는 문제를 역사적·경제적 맥락에 위치시키고, allowance 기반 평판의 동기를 제시하며, 연구 문제·목표·질문을 진술하고, 핵심 구성 개념을 정의하며, 기여, 범위, 학위논문 구성을 개관한다.

\dissertationSubheading[sec:background]{배경: DeFi, 담보화, 가명성}

공개 블록체인은 중앙 운영자 없이 프로그램 가능한 정산과 개방적 참여를 도입하였다 \cend{46,60}. 스마트 계약은 합의 조건을 자체 실행 코드로 인코딩하여, 수탁 중개자가 아닌 온체인에서 정산되는 금융 응용을 가능하게 한다 \cend{9,55}. DeFi 프로토콜은 이러한 계약을 대출, 자동화된 시장 조성, 무기한 선물 시장으로 조합한다 \cend{53,59}. Arbitrum One과 유사한 Layer-2 네트워크에서는 활동 밀도가 충분히 높아, 지갑 수준의 행동 신호---전송, 승인, 프로토콜 이벤트---를 대규모로 관측할 수 있다.

이러한 투명성에도 불구하고 DeFi 신용 시장은 구조적으로 과담보화되어 있다. 프로토콜은 일반적으로 차입자나 거래자가 노출 가치를 초과하는 담보를 예치하도록 요구하는데, 상대방이 식별된 법적 인격이 아니라 가명 주소이기 때문이다 \cend{26,28}. 가명성은 공개 원장의 설계 특성이다. 누구나 무시할 수 있는 비용으로 주소를 생성할 수 있고, 단일 경제 주체가 다수의 지갑을 통제할 수 있다 \cend{17}. 전통적 신용 평가는 신원에 연결된 이력에 의존하며, 이는 본 설정으로 깔끔하게 이전되지 않는다 \cend{48}. 결과적으로 위험 관리는 개별화된 평판보다 담보 완충, 청산 엔진, 프로토콜 수준 매개변수에 의존한다.

담보화는 프로토콜을 보호하지만 자본 비효율을 초래한다. 거래자와 차입자는 다른 곳에 배치할 수 있는 자산을 잠가야 하며, 프로토콜은 확립된 신용을 가진 차입자에게 전통 금융이 제공하는 저담보 상품을 포기한다. 이에 연구자들은 오프체인 신원의 대체물로서 온체인 신호를 탐색해 왔다. 전송 그래프 중심성 \cend{16}, 네트워크 위험 전파 \cend{44}, 전망 이론 및 블록체인 기반 신용 모형 \cend{29}, 대출 프로토콜에 대한 기계학습 파이프라인 \cend{39} 등이 그 예이다. 이들 접근은 공통 전제를 공유한다. 공개 원장 데이터는 사용자를 비익명화하지 않고도 위험 평가에 정보를 제공할 수 있는, 경제적으로 유의미한 관계를 인코딩한다는 것이다.

그래프 순위 방법은 이 전제에 자연스럽게 부합한다. PageRank와 그 변형은 방향성 간선을 통해 영향력을 전파하여, 다른 노드로부터의 재귀적 보증에 따라 노드를 순위화한다 \cend{49,8,40}. EigenTrust와 TrustRank 같은 초기 웹 신뢰 시스템은 유사한 아이디어를 P2P 및 웹 스팸 설정에 적용하였다 \cend{35,27}. Do and Do\cend{16}는 이더리움 전송 그래프에 PageRank와 HodgeRank를 사회적 신용 측정으로 적용한다. Adaptive Weighted PageRank(AWP)\cend{64}는 가치 가중·시간 감쇠 간선을 쓰는 관련 전송 그래프 순위 방법이며, 경제적 중요성에 대한 도메인 직관 대리 지표로서 인바운드 전송량을 기준으로 점수를 검증한다. 이 연구 계보는 전송 그래프 평판을 신뢰할 만한 기준선으로 확립하지만, ERC-20 allowance로 기록되는 명시적 지출 승인---별개의 온체인 행위---은 모형화하지 않는다 \cend{18,19}.

\dissertationSubheading[sec:motivation]{온체인 평판의 동기}

레버리지 DeFi에서 온체인 평판 연구를 동기화하는 네 가지 실무적 압력이 있다.
\begin{itemize}
\item \textbf{자본 효율:} 프로토콜은 온체인 정산을 유지하면서, 예를 들어 한도 상향 주기, 수수료 등급, 모니터링 우선순위와 같은 차별화된 서비스를 지원할 수 있는 서열 신호를 추구하며, 향후 통제된 실험 하에서만 완만한 담보 차별화를 고려한다 \cend{53,26}.
\item \textbf{감사 가능성:} 불투명한 오프체인 신용 점수와 달리, 공개 이벤트로부터 도출된 그래프 기반 점수는 재계산·검사할 수 있어, 블랙박스형 탈중앙화 신용 평점에 대한 우려에 대응한다 \cend{48}.
\item \textbf{운영 비용:} 예를 들어 각 레버리지 조정 전에 자주 갱신해야 하는 평판 시스템은, 매 갱신마다 비용이 큰 과거 전송 집계를 감당할 수 없다.
\item \textbf{표준화:} 주식 시장은 주가수익비율과 주가순자산비율을 완벽한 수익률 예측자가 아니라 스크리닝과 소통을 위한 공유 어휘로 오래 사용해 왔다 \cend{24,6,20}. 가명 DeFi에는 지갑 수준 상대방 차별화를 위한 유사한 서열 언어가 없다 \cend{1}. 제~\ref{ch:discussion}장은 EndorseRank와 AWP가 담보를 대체하거나 거래 이익을 예측한다고 주장하지 않으면서, 수수료·이용률·모니터링 효율을 지원하는 인프라적 역할을 하는 산업 채택 시나리오를 전개한다(제~\ref{sec:practical-implications}절).
\end{itemize}

전송 기반 AWP는 처음 두 압력에는 대응하지만 세 번째에서는 비용이 크다. $G_{\text{transfer}}$를 구축하려면 전송 이력을 집계하고, 로지스틱 시간 감쇠를 적용하며, 밀집한 간선 집합에서 PageRank를 실행해야 한다 \cend{64}. 본 연구에서 사용한 매칭 코호트($n=5{,}521$ 지갑)에서,\footnote{매칭 코호트는 관측 창에서 GMX V2 \texttt{PositionDecrease} 청산이 최소 세 건이고 평판 부분그래프에서 비영(非零) 연결성을 갖는 지갑의 교집합이다. 런타임 확장성은 더 큰 확장 지갑 풀($N=31{,}612$)을 사용한다.} AWP 부분그래프는 $137{,}087$개 간선을 포함하며 $18.711$\,s의 PageRank 벽시계 시간이 필요하다. allowance 그래프는 구조적으로 다른 대안을 제공한다. ERC-20 \texttt{approve} 또는 EIP-2612 \texttt{permit}은 토큰 소유자가 지출자에게 전송을 얼마나 승인하는지를 기록한다 \cend{18,19}. 소유자--지출자 쌍당 최신 allowance는 현재 승인 의도의 스냅샷이다. 신규 승인이 구 승인을 덮어쓰므로 시간 감쇠는 불필요하다. 결과적인 보증 그래프 $G_{\text{approve}}$는 희소하며---동일 코호트에서 $14{,}727$개 간선---PageRank는 $2.105$\,s에 완료되어 AWP보다 약 $8.9\times$ 빠르다. 그 차이는 최신 allowance 스냅샷 하에서의 그래프 희소성($14{,}727$ 대 $137{,}087$ 간선)의 결과이며, 독립된 과학적 기여가 아니다. 주장~C1은 이를 표현의 운영적 속성으로 보고한다.

효율을 넘어 allowance는 구별되는 의미를 지닌다. 전송은 라우팅, 시장 조성, 수동적 수령을 반영할 수 있는 반면, allowance는 지출 권한의 명시적 위임이다. allowance를 보증 간선으로 취급하면 PageRank의 인용 은유 \cend{49}가 프로토콜이 이미 토큰 상호작용에 의존하는 신뢰형 행위와 정렬된다. 그 의미적 차이가 단지 다른 순위를 낳는지, 아니면 경제적으로 유의미한 결과와 다르게 정합하는 순위를 낳는지는 실증적 질문이다. 본 학위논문은 본 체인에서 저담보 대출 부도 라벨이 희소하기 때문에, 밀집한 검증 틀로서 GMX V2 무기한 거래 결과를 사용하여 Arbitrum One에서 그 질문에 답한다 \cend{23}.

\dissertationSubheading[sec:research-problem]{연구 문제}

기존 Adaptive Weighted PageRank(AWP) 방법은 거래 이력에 대한 전송량과 로지스틱 시간 감쇠로부터 신뢰를 추론한다 \cend{64}. 경제적으로 해석 가능하지만, (i) 비용이 큰 과거 전송 집계를 요구하고, (ii) 인바운드 전송량을 ``가치 있는 지갑''의 대리 지표로 취급하며, (iii) 명시적 지출 승인을 모형화하지 않는다. EndorseRank는 대신 소유자--지출자 쌍당 최신 allowance를 방향성 보증 그래프의 간선 강도로 사용하여, 시간 감쇠 가중 없이 더 희소한 그래프와 의미적으로 구별되는 신뢰 신호를 산출한다---즉, 향상된 온체인 지갑 평판 점수화를 위해 ERC-20 allowance 간선을 PageRank에 통합한다.

따라서 연구 문제는 비교·평가적이다. GMX V2 무기한 활동이 있는 동일 Arbitrum One 지갑 표본에서, EndorseRank는 AWP와 구성 개념상 구별되는 지갑 순서를 산출하고, 도메인 직관 대리 지표(allowance 특화 및 거래 성공 계열 포함)와 신뢰할 만하게 정합하며, 현저히 낮은 계산 비용으로 이를 수행하는가? 주요 비교는 EndorseRank 대 AWP이며, 보조적 결과 고유 참조(GF-PR)는 이후 거래 결과 대리 지표의 해석을 한정하기 위해 도입되며, 어느 사회적 방법도 대체하지 않는다.

표~\ref{tab:primary-methods-preview}는 주요 비교를 요약한다. 표~\ref{tab:hrq-claim-map}는 가설(H1--H3), 연구 질문(RQ1--RQ4), 보조 질문 SRQ1, 제4장 실증 주장(C1--C4, SC1)을 요약한다.

\begin{table}[htbp]
\centering
\caption{주요 비교: EndorseRank 대 AWP (매칭 지갑 코호트, $n=5{,}521$).}
\label{tab:primary-methods-preview}
\fitwidth{%
\begin{tabular}{p{0.18\linewidth}p{0.28\linewidth}p{0.48\linewidth}}
\toprule
방법 & 입력 그래프 & 역할 \\
\midrule
EndorseRank & 최신 ERC-20 allowance ($G_{\text{approve}}$) & 주요 구성 개념: allowance 위임; 본 학위논문의 대상 \\
AWP & 시간 감쇠 전송 ($G_{\text{transfer}}$) & 주요 기준선: IEEE RIVF 2023 전송 그래프 방법 \cend{64}; Do and Do\cend{16}는 PageRank/HodgeRank 동반 논문 \\
\bottomrule
\end{tabular}
}
\end{table}

모호성을 피하기 위해, 본 학위논문 전반에서 사용하는 핵심 개념을 다음과 같이 정의한다.

\begin{itemize}
\item \textbf{DeFi (탈중앙화 금융):} 공개 블록체인의 스마트 계약을 통해 구현되는 금융 서비스---대출, 차입, 무기한 거래 포함---로, 일반적으로 수탁이나 정산을 통제하는 중앙 운영자가 없다 \cend{53,28}. DeFi 프로토콜은 개방 인터페이스를 조합하여 포지션, 청산, 토큰 이동이 온체인에서 관측 가능하게 한다.
\item \textbf{스마트 계약:} 신뢰할 수 있는 중개자 없이 합의 조건을 강제하는, 블록체인에 배포된 자체 실행 프로그램 코드 \cend{9,55}. 본 학위논문에서 스마트 계약은 ERC-20 allowance 로그, 전송 로그, GMX V2 포지션 이벤트의 출처이다.
\item \textbf{지갑 / 주소:} 자산을 보유하고 스마트 계약과 상호작용하는 블록체인 계정; 평판 점수화의 분석 단위 \cend{9,18}. 지갑은 주소로 식별되며, 오프체인 신원은 가정하거나 복원하지 않는다.
\item \textbf{온체인 평판 점수:} 가명성 하에서 금융 상호작용을 위한 지갑의 상대적 위치를 나타내는, 온체인 신호로부터 도출된 수치 지표 \cend{16}. 점수는 순위화 목적의 서열이며, 절대 크기는 신용 한도로 해석하지 않는다.
\item \textbf{온체인 신용 위험:} 공개 블록체인에서 포지션 청산이나 반복적 거래 손실과 같은 관측 가능한 불리한 결과 \cend{44}. 본 학위논문은 GMX V2 무기한 결과를 위험 관련 행동에 대한 도메인 직관 대리 지표로 사용하며, 대출 상환 정답(ground truth)으로 사용하지 않는다.
\item \textbf{GMX V2 무기한 스왑 포지션:} Arbitrum One에서 \textbf{PositionIncrease}로 개설되고 \textbf{PositionDecrease}로 청산되는, GMX V2의 담보화된 롱 또는 숏 노출 \cend{23,53}. 청산으로부터의 실현 손익이 거래 성공 및 역위험 대리 지표를 공급한다.
\item \textbf{PageRank:} 감쇠 계수 하에서 방향성 간선을 통해 영향력을 전파하는 그래프 순위 방법 \cend{49,8,40}. \textbf{EndorseRank}는 최신 allowance 보증 그래프에 PageRank를 적용한다. 간선 강도는 소유자--지출자 쌍당 현재 allowance와 같으며 시간 감쇠는 없다. \textbf{Adaptive Weighted PageRank(AWP)}는 전송 가치와 로지스틱 시간 감쇠로 주어진 간선 가중치를 갖는 전송 그래프에 PageRank를 적용한다 \cend{64}.
\item \textbf{ERC-20 allowance / Approve / Permit:} 토큰 소유자가 지출자에게 소유자를 대신하여 전송을 얼마나 허용하는지를 기록하는 메커니즘 \cend{18,19}. \texttt{Approve}는 표준 온체인 호출이며, EIP-2612 \texttt{permit}은 서명을 통해 동일한 상태 변경을 승인한다. EndorseRank는 소유자--지출자 쌍당 최신 비영(非零) allowance 스냅샷을 사용한다.
\item \textbf{보증 그래프:} 단순 자산 전송이 아니라 지출 권한의 명시적 위임을 간선이 나타내는 방향성 그래프 \cend{18,19}. EndorseRank에서 간선 $u \to v$는 지갑 $u$가 현재 지갑 또는 계약 $v$에게 $u$가 소유한 토큰을 지출하도록 승인함을 의미한다. 전송 그래프 AWP는 대비되는 기준선을 제공한다 \cend{16}.
\item \textbf{도메인 정합:} 평판 순위가 관측 가능한 대리 지표 순위와 일치하는 정도이며, Spearman $\rho$와 Kendall $\tau$로 평가한다 \cend{54,36,14,16}. 정합은 구성 개념별이다. 한 방법은 자신의 주요 그래프 구성 개념과 강하게 정합하고 직교하는 결과 계열과는 약하게 정합할 수 있다.
\end{itemize}

\dissertationSubheading[sec:supplementary-extension]{보조 확장 (GF-PR)}

실증 분석 과정에서 \textbf{GF-PR}(GainFlow PageRank)를 보조적 결과 고유 참조로 추가하였다(제~\ref{sub:gf-pr}절). GF-PR는 검증 대리 지표를 정의하는 데 사용한 동일 GMX 손익 흐름으로부터 방향성 스타 그래프를 구축한다. 사회적 지갑 간 간선이 아니라 합성 POOL/SINK 흐름으로 지갑을 순위화하여, 진단적 해석을 위한 거래 결과 대리 지표 정합의 상한을 설정한다---배포 가능한 사회적 평판 점수가 아니다.

\begin{table}[htbp]
\centering
\caption{실증 과정에서 추가된 보조 확장 (주요 EndorseRank--AWP 연구 질문의 일부가 아님).}
\label{tab:supplementary-gf-pr}
\fitwidth{%
\begin{tabular}{p{0.18\linewidth}p{0.28\linewidth}p{0.48\linewidth}}
\toprule
방법 & 입력 그래프 & 역할 \\
\midrule
GF-PR & GMX PnL 스타 (POOL$\to$지갑, 지갑$\to$SINK) & 보조 상한 참조; 거래 결과 대리 지표에서의 구성 개념 중첩이 예상됨 \\
\bottomrule
\end{tabular}
}
\end{table}

\begin{table}[htbp]
\centering
\caption{가설, 연구 질문, 실증 주장 요약.}
\label{tab:hrq-claim-map}
\fitwidth{%
\begin{tabular}{p{0.12\linewidth}p{0.42\linewidth}p{0.40\linewidth}}
\toprule
항목 & 다루는 내용 & 보고 위치 \\
\midrule
H1 / RQ1 & EndorseRank와 AWP가 구별되는 지갑 순서를 산출하는가? & \S\ref{sec:rank-divergence}; 방법 간 $\tau=0.316$ \\
H2 / RQ2 & 어느 사회적 방법이 어느 대리 지표 계열과 더 강하게 정합하는가? & 표~\ref{tab:alignment-three-methods}--\ref{tab:alignment-gmx}; 주장 C2--C3 \\
H3 / RQ3 & 동일 $n$에서 EndorseRank가 AWP보다 빠른가? & 표~\ref{tab:benchmark-runtime}--\ref{tab:benchmark-scaling}; 주장 C1, C4 \\
RQ4 & 구성 타당도 관측: 전송 대 거래 성공 틀 & 표~\ref{tab:alignment-three-methods} \\
SRQ1 & GF-PR 상한이 거래 결과 대리 지표 해석에 함의하는 바는? & \S\ref{sec:gf-pr-ceiling}; 주장 SC1 \\
C1 & EndorseRank 대 AWP 런타임과 그래프 크기 & 표~\ref{tab:benchmark-runtime} \\
C2 & 선택된 대리 지표에서 신뢰할 만한 EndorseRank--AWP 정합 & 표~\ref{tab:alignment-transfer}--\ref{tab:alignment-gmx} \\
C3 & 구성 개념별 우세 (EndorseRank 대 AWP) & 표~\ref{tab:alignment-three-methods} \\
C4 & EndorseRank의 효율--정합 상충 & 표~\ref{tab:benchmark-runtime}--\ref{tab:benchmark-scaling} \\
SC1 & 보조 GF-PR 구성 개념 중첩 진단 & \S\ref{sec:gf-pr-ceiling} \\
\bottomrule
\end{tabular}
}
\end{table}

\dissertationSubheading[sec:research-objectives]{연구 목표}

본 연구의 주요 목표는 RQ1--RQ4에 일대일로 대응한다.

\begin{enumerate}
\item[\textbf{O1.}] 동일 Arbitrum One 코호트에서 EndorseRank와 AWP가 구별되는 지갑 순서를 유도하는지를 확인한다(RQ1).
\item[\textbf{O2.}] 다섯 대리 지표 계열에 대한 각 사회적 방법의 구성 개념별 정합을 측정한다(RQ2).
\item[\textbf{O3.}] 동일 $n$에서 EndorseRank와 AWP의 PageRank 계산 비용을 비교하고, 확장 풀 확장성을 포함한다(RQ3).
\item[\textbf{O4.}] 본 체인에서 사회적 방법에 대해 어느 검증 틀---전송 중심성 또는 거래 성공---이 더 강한 순위 정합을 산출하는지를 관측한다(RQ4).
\end{enumerate}

\dissertationSubheading[sec:research-questions]{연구 질문}

본 연구는 세 가지 주요 가설을 검정한다. \textbf{H1}---EndorseRank와 AWP는 동일 코호트에서 구별되는 지갑 순서를 산출한다. \textbf{H2}---두 사회적 방법은 구성 개념별 대리 지표 정합에서 다르다. \textbf{H3}---동일 $n$에서 EndorseRank가 AWP보다 계산적으로 더 효율적이다. 다음 연구 질문이 표~\ref{tab:hrq-claim-map}에 요약된 실증 평가를 안내한다. RQ1--RQ4는 주요 EndorseRank--AWP 비교를 구성하며, SRQ1은 보조 GF-PR 확장을 다룬다.

\begin{enumerate}
\item[\textbf{RQ1.}] EndorseRank와 AWP는 동일 Arbitrum 지갑 표본에서 의미 있게 다른 지갑 순서를 산출하는가?
\item[\textbf{RQ2.}] EndorseRank와 AWP 중, 어느 사회적 방법이 다섯 대리 지표 계열 각각과 더 강하게 정합하는가(Spearman $\rho$, Kendall $\tau$)?
\item[\textbf{RQ3.}] 동일 매칭 표본에서 EndorseRank가 AWP보다 계산적으로 더 효율적인가(런타임, 메모리, 반복 횟수)?
\item[\textbf{RQ4.}] 우열 경기가 아니라 구성 타당도 관측으로서: 본 체인에서 전송 대리 지표(진입차수/진입가치)와 거래 성공 대리 지표(GMX V2 무기한 청산) 중 어느 쪽이 사회적 평판 방법에 더 강한 순위 정합을 산출하는가?
\item[\textbf{SRQ1.}] 보조 GF-PR 결과 고유 상한은 EndorseRank와 AWP의 거래 결과 대리 지표 정합을 해석하는 데 무엇을 함의하는가?
\end{enumerate}

RQ1은 allowance 그래프와 전송 그래프가 구별되는 중심성 구조를 유도하는지를 검정하며, 점수 벡터 간 방법 간 Kendall $\tau$로 조작화한다. RQ2는 전송, allowance, 역위험, Sybil 조정 안정성, 거래 성공 대리 지표 계열에 걸쳐 구성 개념별 정합을 평가한다. RQ3는 동일 지갑 집합에서 PageRank 벽시계 시간, 메모리, 반복 횟수를 비교하며, 확장 지갑 풀에서의 확장성을 포함한다. RQ4는 본 표본에서 사회적 그래프 점수에 어느 외부 검사가 더 정보적인지에 대한 관측이며, 한 틀이 평판의 보편적 정의라고 주장하지 않는다. SRQ1은 결과 흐름에 직접 순위를 매길 때 거래 결과 정합이 얼마나 높아질 수 있는지의 상한을 설정하기 위해 GF-PR만을 사용하여, 구성 개념 중첩을 명확히 한다 \cend{14}.

\dissertationSubheading[sec:contributions]{기여}

본 학위논문은 온체인 평판 연구와 DeFi 위험 평가에 네 가지 기여를 한다.
\begin{itemize}
\item \textbf{EndorseRank:} 최신 ERC-20 \texttt{approve}/\texttt{permit} 간선에 PageRank를 적용하는 allowance 기반 보증 그래프 모형. EndorseRank는 새로운 순위 알고리즘이 아니며, 고전적 PageRank\cend{49}를 새로운 간선 구성 개념(최신 ERC-20 allowance)에 적용한다. 현재 지출 승인을 보증 강도로 취급하고, allowance 덮어쓰기 의미론에 의존하여 시간 감쇠를 제거하며, 빈번한 갱신에 적합한 희소 그래프를 산출한다.
\item \textbf{5개 대리 지표 도메인 정합 평가:} 레버리지 DeFi 활동에서 EndorseRank와 AWP를 비교한다. 대리 지표는 전송 중심성 \cend{16}, allowance 위임, GMX V2 역위험 및 거래 성공 \cend{23}, Sybil 조정 안정성을 포괄한다. 정합은 Spearman $\rho$와 Kendall $\tau$로 보고되어 \cend{54,36}, 단일 전역 우세자가 아닌 구성 개념별 해석을 가능하게 한다.
\item \textbf{개방형 재현 가능 평가 파이프라인:} BigQuery 추출, 그래프 구축, PageRank 점수화, 대리 지표 계산, LaTeX 결과 표를 기계 판독 가능 요약에 연결한다(부록~\ref{ch:appendix-eval}). 파이프라인은 주장 C1--C4 및 SC1의 독립적 검증을 지원한다.
\item \textbf{GF-PR 상한 분석:} 매칭 지갑 코호트($n=5{,}521$)에서 거래 결과 대리 지표의 구성 개념 중첩을 명확히 한다. 실증 과정에서 보조 기여로 추가되었으며, GF-PR는 배포 가능한 사회적 평판 방법으로 제안되지 않고, 결과 고유 계열에서 높은 $\tau$의 해석을 한정한다(주장~SC1).
\end{itemize}

방법 간 $\tau$, 계열 수준 정합, 런타임의 핵심 수치는 제~\ref{ch:results}장(표~\ref{tab:benchmark-runtime}--\ref{tab:alignment-three-methods})에서 보고하며 여기서 반복하지 않는다. 그 결과는 어느 사회적 방법의 보편적 우세가 아니라 효율--정합 상충을 지지한다.

\dissertationSubheading[sec:scope-limitations]{범위와 한계}

본 학위논문은 관측 창 2025-12-01부터 2026-05-31까지 Arbitrum One과 GMX V2 무기한 거래에 초점을 둔다. 주요 실증 표본은 GMX \texttt{PositionDecrease} 이벤트가 최소 세 건이고 비영(非零) 평판 부분그래프 연결성을 갖는 $n=5{,}521$ 주소의 매칭 지갑 코호트이며, $365{,}488$건의 디코딩된 청산, $14{,}727$개 EndorseRank 간선, $137{,}087$개 AWP 간선으로 구성된다. 런타임 확장성은 확장 지갑 풀($N=31{,}612$)에서 보고한다. 강건성 검증은 감쇠 민감도, 상위 토큰 부분그래프, 표본 크기 민감도를 검토한다.

무담보 대출 부도 라벨은 범위 밖이다. 저담보 신용 사건은 본 체인에서 신뢰할 만한 정답 검증에 너무 희소하다. 구성 개념과 윤리는 제~\ref{ch:introduction}장과 제~\ref{ch:methodology}장에서 정의한다. 사회적 방법 가운데, 본 표본에서는 전송 대리 지표가 거래 성공 대리 지표보다 더 강하게 정합한다(RQ4). 이는 EndorseRank의 allowance 특화 타당성이나 런타임 우위를 약화시키지 않는다. GF-PR 결과는 보조 진단으로 보고되며(SRQ1), 세 번째 대등 평판 방법으로 보고되지 않는다. 향후 비교를 위해 보관된 확장 기준선은 본문 서술에서 제외한다(부록~\ref{ch:appendix-eval}).

Arbitrum One, GMX V2, 6개월 창을 넘는 일반화는 주장하지 않는다. 토큰 소수점 정규화, 가격 오라클, permit 대 approve 커버리지는 간선 가중치에 영향을 줄 수 있다. allowance나 전송의 적대적 조작은 제~\ref{ch:discussion}장에서 논의하는 위협으로 남으며, Sybil 조정 대리 지표는 그 위험을 부분적으로 다루지만 제거하지는 않는다 \cend{17}.

\dissertationSubheading[sec:roadmap]{학위논문 구성}

학위논문은 다음과 같이 진행된다.

제~\ref{ch:literature}장은 그래프 기반 온체인 평판, ERC-20 allowance 의미론, GMX V2 거래 성공 검증을 검토한다. 주요 EndorseRank 대 AWP 비교를 동기화하는 이론적 틀---하이퍼링크 전파, 도메인 정합, 계산 복잡도---을 전개하고, GF-PR를 보조적 결과 고유 상한으로 위치시킨다.

제~\ref{ch:methodology}장은 데이터 출처, 표본추출, EndorseRank와 AWP의 그래프 구축, 5개 대리 지표 평가 설계, 계산 벤치마크, 강건성 검증, 재현성 관행, 윤리를 명세한다. 매칭 지갑 코호트, 확장성에 사용하는 확장 지갑 풀, 제~\ref{ch:results}장 표를 산출하는 조작 절차를 정의한다.

제~\ref{ch:results}장은 실증 결과를 보고한다. 코호트 특성, 런타임 및 확장성 벤치마크(주장~C1), 강건성 검증, 다섯 대리 지표 계열에 걸친 사회적 방법 정합(주장~C2--C3), 순위 분기(H1/RQ1), 보조 GF-PR 상한 분석(SRQ1/SC1). 주요 주장 C1--C4와 보조 주장 SC1을 명시적으로 진술한다.

제~\ref{ch:discussion}장은 연구 질문별로 발견을 해석하고, 이론적·실무적 함의를 전개하며, allowance 그래프에서의 Sybil 공격 비용 모형을 정식화하고, 규제·윤리적 고려를 논의하며, 타당성에 대한 위협과 한계를 평가한다.

제~\ref{ch:conclusion}장은 지식에 대한 기여를 종합하고, RQ1--RQ4 및 SRQ1에 대한 발견을 요약하며, 교차 프로토콜 검증, 적대적 강건성, 하이브리드 방법, 실시간 오라클 통합에 대한 향후 연구를 개관한다.
